% STRATA: JMLR Manuscript
\documentclass[twoside,11pt]{article}

\usepackage[preprint]{jmlr2e}
\usepackage{amsmath,amsfonts,amssymb}
\usepackage{graphicx}
\usepackage{booktabs}
\usepackage{multirow}
\usepackage{algorithm}
\usepackage{algorithmic}
\usepackage{lastpage}

\newtheorem{assumption}{Assumption}

\jmlrheading{0}{2026}{1-\pageref{LastPage}}{3/26; Revised 3/27}{3/27}{0000}{Matthew Powell}
\ShortHeadings{STRATA for Heterogeneous Conformal Prediction}{Powell}
\firstpageno{1}

\begin{document}

\title{STRATA: Conformal Prediction with Heterogeneous Message-Passing Calibration for Infrastructure Risk Assessment}

\author{\name Matthew Powell \email matthew.a.powell@outlook.com \\
        \addr Independent Researcher\\
        GIS and Machine Learning Research\\
        United States}

\editor{To be assigned}

\maketitle

% ─────────────────────────────────────────────────────────────
\begin{abstract}%
We study conformal prediction for heterogeneous graphs where node difficulty depends on both graph topology and node type. STRATA combines a heterogeneous graph neural network with Conformal Heterogeneous Message Passing (CHMP), which rescales conformal scores using frozen training residuals aggregated over typed neighborhoods. This preserves split-conformal validity while adapting interval width to local graph difficulty. Across 20 synthetic seeds, all calibrators maintain near-nominal 90\% coverage, while an ensemble variant yields the narrowest intervals ($0.743 \pm 0.052$). On two public MATPOWER-based benchmarks, ACTIVSg200 and IEEE 118, STRATA attains average marginal coverage of $0.935$ and $0.898$, respectively. Relative to a Mondrian baseline on these real-data runs, CHMP matches average coverage, produces a small width reduction on ACTIVSg200, and shows no consistent calibration gain on IEEE 118. Additional ablations show the expected width-coverage tradeoff as $\alpha$ varies, and a group-wise audit yields a worst-group coverage gap of $0.035$ in a representative run. These results position STRATA as a general methodology for uncertainty-aware prediction on heterogeneous graph-structured data while clarifying that CHMP's empirical gains are strongest in the synthetic benchmark.
\end{abstract}

\begin{keywords}
conformal prediction, graph neural networks, heterogeneous graphs, uncertainty quantification, calibration
\end{keywords}

% ─────────────────────────────────────────────────────────────
\section{Introduction}
\label{sec:intro}

Critical infrastructure systems---electric power, water distribution,
and telecommunications---form interdependent networks where failures
cascade across utility boundaries
\citep{rinaldi2001identifying,buldyrev2010catastrophic}.
A transformer outage at a power substation can disable water treatment
pumps, which in turn disrupts telecom cooling systems at co-located
facilities \citep{ouyang2014review}. Predicting node-level risk scores
in such coupled systems requires not only accurate point predictions
but also reliable uncertainty quantification: operators need to know
\\emph{how wrong} a risk estimate might be at each infrastructure
node, and that uncertainty is inherently non-uniform across the
heterogeneous graph.

\textbf{Conformal prediction} \citep{vovk2005algorithmic,shafer2008tutorial,
angelopoulos2023conformal} offers distribution-free prediction intervals with
finite-sample coverage guarantees, making it attractive for safety-critical
infrastructure applications where parametric assumptions may be violated.
However, standard split conformal prediction produces intervals of uniform
width across all nodes, ignoring the topology-dependent difficulty structure
that arises from infrastructure coupling.

\textbf{Graph neural networks} (GNNs) \citep{kipf2017semi,gilmer2017neural,
hamilton2017inductive} have emerged as powerful tools for node-level prediction
on relational data, and recent work has extended conformal prediction to
graph-structured settings \citep{zargarbashi2023conformal,huang2024uncertainty,
clarkson2023distribution}. However, all existing graph conformal prediction
methods target \emph{node classification}---producing prediction
\emph{sets}---and operate on homogeneous, single-type graphs.
To date, no framework provides distribution-free prediction
\emph{intervals} for node-level \emph{regression} on heterogeneous
multi-type graphs, nor does any prior method exploit graph-propagated
residuals to adapt interval widths while preserving finite-sample coverage
guarantees.

We address these gaps with \textbf{STRATA}, the first conformal prediction
framework for heterogeneous graph regression, integrating
three key innovations:

\begin{enumerate}
\item \textbf{Conformal Heterogeneous Message Passing (CHMP)}: A
propagation-aware calibration scheme that normalizes nonconformity scores using
frozen training-set residuals propagated through the infrastructure coupling
topology. Rather than reducing aggregate interval width, CHMP
\emph{redistributes} width across nodes---giving wider intervals to nodes in
high-difficulty neighborhoods and narrower intervals to low-difficulty nodes.

\item \textbf{A family of advanced calibrators} that learn the normalization
function from data: a MetaCalibrator using heteroscedastic Gaussian NLL
\citep{kendall2017uncertainties}, an AttentionCalibrator using learned neighbor
weights, a LearnableLambdaCalibrator with per-type grid-search, and a
heterogeneous CQR variant \citep{romano2019conformalized}. An ensemble-based
calibrator using epistemic variance \citep{lakshminarayanan2017simple} achieves
the strongest width reduction ($6.4\%$ narrower than Mondrian CP).

\item \textbf{Spatial diagnostics} including Moran's I autocorrelation tests
\citep{moran1950notes}, Getis-Ord $G_i^*$ hotspot detection
\citep{getis1992analysis}, and conformal kriging surfaces \citep{cressie1993statistics}
that bridge infrastructure risk assessment with geostatistical analysis.
\end{enumerate}

% ─────────────────────────────────────────────────────────────
\section{Related Work}
\label{sec:related}

\subsection{Conformal Prediction}

Conformal prediction, developed by \citet{vovk2005algorithmic}, provides
distribution-free prediction regions with finite-sample coverage guarantees under
the exchangeability assumption. The inductive (split) conformal framework
\citep{papadopoulos2002inductive} partitions data into training, calibration,
and test sets. \citet{lei2018distribution} established distribution-free
inference guarantees for regression. Mondrian conformal prediction
\citep{vovk2012conditional} extends coverage guarantees to group-conditional
settings.

Recent extensions have relaxed the exchangeability requirement.
\citet{barber2023conformal} proved finite-sample coverage results beyond
exchangeability, directly motivating our use of normalized scores on
graph-structured data. \citet{gibbs2021adaptive} proposed adaptive conformal
inference for streaming settings with distribution shift.
Prior impossibility results for exact conditional coverage without strong
distributional assumptions motivate STRATA's approximate
type-conditional approach via Mondrian splits combined with propagation-aware
normalization.

\subsection{Conformal Prediction on Graphs}

\citet{zargarbashi2023conformal} introduced DAPS, a diffusion-based approach for
conformal prediction \emph{sets} on GNNs for node classification.
\citet{clarkson2023distribution} provided distribution-free coverage guarantees
for node classification under graph dependence, and
\citet{huang2024uncertainty} developed CF-GNN for node-level uncertainty
quantification, also in the classification setting.
Most recently, \citet{song2024snaps} proposed SNAPS, using
similarity-navigated scores for graph conformal classification (NeurIPS'24).
Notably, \textbf{all prior graph conformal prediction work targets
classification}---producing prediction \emph{sets}---and operates on
homogeneous, single-type graphs.
STRATA differs by (i)~targeting \emph{regression} with prediction
\emph{intervals}, (ii)~modeling heterogeneous multi-typed infrastructure
graphs with per-type Mondrian coverage, and (iii)~introducing
propagation-aware normalization (CHMP) using frozen training-set residuals.

\subsection{Graph Neural Networks for Heterogeneous Data}

Heterogeneous GNNs extend message passing to multi-typed graphs: R-GCN
\citep{schlichtkrull2018modeling} introduces relation-specific weight matrices,
HAN \citep{wang2019heterogeneous} applies hierarchical attention, and HGT
\citep{hu2020heterogeneous} uses transformer-style attention. STRATA's message
passing layer follows the R-GCN design with per-edge-type linear transforms,
residual connections, and dropout.

\subsection{Uncertainty Quantification}

Deep ensembles \citep{lakshminarayanan2017simple} provide calibrated epistemic
uncertainty. \citet{kendall2017uncertainties} decomposed uncertainty into
aleatoric and epistemic components via heteroscedastic loss---directly inspiring
STRATA's MetaCalibrator. Conformalized quantile regression (CQR)
\citep{romano2019conformalized} combines neural network quantile regression
\citep{koenker1978regression} with conformal calibration. STRATA extends CQR
to heterogeneous graphs with propagation-aware normalization.

% ─────────────────────────────────────────────────────────────
\section{Methods}
\label{sec:methods}

\subsection{Problem Formulation}

Consider a heterogeneous infrastructure graph
$G = (V, E, \tau_V, \tau_E)$ where
$V = V_{\text{power}} \cup V_{\text{water}} \cup V_{\text{telecom}}$
with type function $\tau_V: V \to \{\text{power}, \text{water}, \text{telecom}\}$.
Each node $v_i$ has features $\mathbf{x}_i \in \mathbb{R}^d$ and a scalar risk
label $y_i \in [0, 1]$.

Given a trained GNN with point predictions $\hat{y}_i$, we seek prediction
intervals $C_i = [\ell_i, u_i]$ such that:
\begin{equation}
\Pr(y_i \in C_i) \geq 1 - \alpha
\end{equation}
for a user-specified miscoverage level $\alpha$, ideally with the additional
type-conditional guarantee:
\begin{equation}
\Pr(y_i \in C_i \mid \tau_V(i) = t) \geq 1 - \alpha, \quad
\forall t \in \{\text{power}, \text{water}, \text{telecom}\}
\end{equation}

\subsection{Heterogeneous GNN Architecture}

STRATA's GNN backbone extends the R-GCN framework \citep{schlichtkrull2018modeling}
with residual connections. Each \texttt{HeteroMessagePassingLayer} $\ell$ computes:
\begin{equation}
\mathbf{h}_i^{(\ell)} = \text{ReLU}\!\left(\mathbf{W}_{\text{self}}
\mathbf{h}_i^{(\ell-1)} + \sum_{r \in \mathcal{R}}
\frac{1}{|\mathcal{N}_r(i)|} \sum_{j \in \mathcal{N}_r(i)}
\mathbf{W}_r \mathbf{h}_j^{(\ell-1)}\right) + \mathbf{h}_i^{(\ell-1)}
\end{equation}
where $\mathcal{R}$ is the set of edge types and $\mathbf{W}_r$ are
per-relation weight matrices. Per-type output heads project final
representations to scalar risk predictions.

\subsection{Conformal Heterogeneous Message Passing (CHMP)}

CHMP introduces propagation-aware normalization:

\textbf{Step 1.} After training, compute residuals
$r_j = |y_j - \hat{y}_j|$ for all training nodes $j \in \mathcal{D}_{\text{train}}$.

\textbf{Step 2.} Aggregate neighbor difficulty:
\begin{equation}
\bar{r}_{\mathcal{N}(i)} = \text{agg}_{j \in \mathcal{N}(i) \cap
\mathcal{D}_{\text{train}}} (r_j)
\end{equation}

\textbf{Step 3.} Compute normalization signal:
$\sigma_i = 1 + \lambda \cdot \bar{r}_{\mathcal{N}(i)}$

\textbf{Step 4.} Calibration scores: $s_i = |y_i - \hat{y}_i| / \sigma_i$

\textbf{Step 5.} Prediction intervals:
\begin{equation}
C_i = [\hat{y}_i - \hat{q}_{\tau(i)} \cdot \sigma_i, \;
\hat{y}_i + \hat{q}_{\tau(i)} \cdot \sigma_i]
\end{equation}

\begin{theorem}[CHMP Coverage Guarantee]
\label{thm:chmp}
Let $\{(X_i, Y_i)\}_{i=1}^{n+1}$ be exchangeable random variables, let
$\hat{f}$ be trained on $\mathcal{D}_{\text{train}}$ (disjoint from
calibration and test sets), and let $\sigma_i = 1 + \lambda \cdot
\bar{r}_{\mathcal{N}(i)}$ where $\bar{r}_{\mathcal{N}(i)}$ is computed
exclusively from frozen training-set residuals. Define normalized scores
$s_i = |Y_i - \hat{f}(X_i)| / \sigma_i$ for
$i \in \mathcal{D}_{\text{cal}} \cup \{n+1\}$. Let $\hat{q}$ be the
$\lceil(1-\alpha)(|\mathcal{D}_{\text{cal}}|+1)\rceil /
|\mathcal{D}_{\text{cal}}|$ quantile. Then:
\begin{equation}
\Pr\!\left(Y_{n+1} \in [\hat{f}(X_{n+1}) - \hat{q} \cdot \sigma_{n+1},
\; \hat{f}(X_{n+1}) + \hat{q} \cdot \sigma_{n+1}]\right) \geq 1 - \alpha
\end{equation}
\end{theorem}

\begin{proof}[Proof sketch]
Since $\sigma_i$ is a deterministic function of training data, the normalized
scores $s_i$ inherit exchangeability from the original variables
\citep{barber2023conformal}. The result follows from the standard split
conformal guarantee \citep{papadopoulos2002inductive,lei2018distribution}.
The Mondrian extension to per-type quantiles similarly preserves
type-conditional coverage \citep{vovk2012conditional}.
\end{proof}

\begin{assumption}
(A1) Exchangeability of calibration and test data conditioned on the training
set. (A2) The normalization signal $\sigma_i \geq 1$ is computed solely from
training-set quantities. (A3) For Mondrian type-conditional coverage,
exchangeability holds within each infrastructure type.
\end{assumption}

\subsection{Advanced Calibrator Variants}

\paragraph{MetaCalibrator.} A lightweight MLP predicts $\sigma_i$ from
node-level features including the point prediction, neighbor residual
statistics, and node degree, trained with heteroscedastic Gaussian NLL.

\paragraph{AttentionCalibrator.} Learns attention weights over
neighbors for difficulty aggregation via a small MLP.

\paragraph{LearnableLambdaCalibrator.} Optimizes $\lambda$ per type through a
50/50 calibration-tune split with grid search.

\paragraph{CQR with Propagation.} Extends conformalized quantile regression
\citep{romano2019conformalized} to heterogeneous graphs with optional
$\sigma_i$ normalization. Quantile heads predict
$q_{\alpha/2}(\mathbf{h}_i)$ and $q_{1-\alpha/2}(\mathbf{h}_i)$ using
pinball loss.

\subsection{Ensemble Uncertainty Decomposition}

$M$ independent GNNs with different random seeds provide:
$\hat{y}_i = \frac{1}{M}\sum_m \hat{y}_i^{(m)}$,
$\text{Var}_i = \frac{1}{M}\sum_m (\hat{y}_i^{(m)} - \hat{y}_i)^2$.
The EnsembleCalibrator uses $\sigma_i = 1 + \lambda \sqrt{\text{Var}_i}$.

% ─────────────────────────────────────────────────────────────
\section{Experimental Setup}
\label{sec:experiments}

\subsection{Synthetic Infrastructure Data}

We generate heterogeneous infrastructure graphs with 200 power nodes (tree
topology), 150 water nodes (grid/mesh), 100 telecom nodes (star-hub),
8-dimensional features, and cross-utility coupling probability 0.3.

\subsection{Real Infrastructure Data: ACTIVSg200}

We use the ACTIVSg200 200-bus Central Illinois power grid
\citep{birchfield2017grid}, constructing a three-layer heterogeneous
graph with power (200 nodes), water ($\sim$108 nodes derived from demand
centers), and telecom ($\sim$80 hub nodes).

\subsection{Additional Public Benchmark: IEEE 118}

To reduce the risk that conclusions depend on a single real-network topology, we
also evaluate STRATA on the public IEEE 118-bus MATPOWER case. As with
ACTIVSg200, we retain the real transmission-layer graph and derive water and
telecom layers from demand centres, yielding a second public heterogeneous
benchmark with a distinct topology and degree distribution.

\subsection{Evaluation Protocol}

All experiments use 20 random seeds with 60/20/20 train/cal/test splits.
Metrics: marginal coverage, type-conditional coverage, mean interval width,
ECE \citep{guo2017calibration}, paired Wilcoxon tests
\citep{wilcoxon1945individual}, and Friedman test \citep{friedman1937use}.

% ─────────────────────────────────────────────────────────────
\section{Results}
\label{sec:results}

\subsection{Baseline Comparison}

\begin{table}[ht]
\centering
\caption{Baseline comparison (20 seeds, $\alpha = 0.10$).}
\label{tab:baseline}
\begin{tabular}{lccc}
\toprule
Method & Coverage & Width & ECE \\
\midrule
Mondrian CP & $0.908 \pm 0.032$ & $0.794 \pm 0.072$ & $0.057 \pm 0.022$ \\
CHMP (mean) & $0.908 \pm 0.030$ & $0.793 \pm 0.071$ & $0.080 \pm 0.020$ \\
CHMP (median) & $0.908 \pm 0.030$ & $0.794 \pm 0.071$ & $0.078 \pm 0.021$ \\
CHMP (med.+floor) & $0.908 \pm 0.032$ & $0.794 \pm 0.072$ & $0.067 \pm 0.023$ \\
\bottomrule
\end{tabular}
\end{table}

All methods achieve the nominal 90\% coverage target. Coverage is
statistically indistinguishable (Friedman $p = 0.93$), confirming CHMP
preserves the conformal guarantee. CHMP's primary effect is width
\emph{redistribution} rather than aggregate reduction.

\subsection{Advanced Calibrator Comparison}

\begin{table}[ht]
\centering
\caption{Full method comparison (20 seeds, $\alpha = 0.10$).
Best width in \textbf{bold}.}
\label{tab:full}
\begin{tabular}{lccc}
\toprule
Method & Coverage & Width & ECE \\
\midrule
Mondrian CP & $0.908 \pm 0.032$ & $0.794 \pm 0.072$ & $\mathbf{0.057}$ \\
CHMP (mean) & $0.908 \pm 0.030$ & $0.793 \pm 0.071$ & $0.080$ \\
CHMP (median) & $0.908 \pm 0.030$ & $0.794 \pm 0.071$ & $0.078$ \\
CHMP (med.+floor) & $0.908 \pm 0.032$ & $0.794 \pm 0.072$ & $0.067$ \\
MetaCalibrator & $0.916 \pm 0.027$ & $0.792 \pm 0.069$ & $0.075$ \\
AttentionCalibr. & $0.914 \pm 0.027$ & $0.790 \pm 0.070$ & $0.070$ \\
LearnableLambda & $0.914 \pm 0.026$ & $0.788 \pm 0.064$ & $0.065$ \\
CQR + propagation & $0.900 \pm 0.034$ & $0.838 \pm 0.063$ & $0.083$ \\
\textbf{Ensemble (M=3)} & $0.920 \pm 0.030$ & $\mathbf{0.743 \pm 0.052}$ & $0.073$ \\
\bottomrule
\end{tabular}
\end{table}

The ensemble-based calibrator achieves the narrowest intervals ($0.743$), a
$6.4\%$ reduction from Mondrian CP, while maintaining valid coverage ($0.920$).
CQR with propagation produces the widest intervals ($0.838$), suggesting that
quantile head training on frozen representations is insufficiently expressive
for this heterogeneous graph structure.

\subsection{Lambda Sensitivity}

\begin{table}[ht]
\centering
\caption{Lambda sensitivity (20 seeds, $\alpha = 0.10$).}
\label{tab:lambda}
\begin{tabular}{cccc}
\toprule
$\lambda$ & Coverage & Width & ECE \\
\midrule
0.0 & $0.916 \pm 0.033$ & $0.792 \pm 0.078$ & $0.056$ \\
0.1 & $0.916 \pm 0.032$ & $0.791 \pm 0.077$ & $0.069$ \\
0.3 & $0.916 \pm 0.029$ & $0.789 \pm 0.075$ & $0.067$ \\
0.5 & $0.915 \pm 0.029$ & $0.788 \pm 0.073$ & $0.069$ \\
1.0 & $0.912 \pm 0.029$ & $0.791 \pm 0.071$ & $0.071$ \\
\bottomrule
\end{tabular}
\end{table}

Coverage remains stable across all $\lambda$ values (within $0.4\%$). Width
shows minimal sensitivity ($<0.5\%$ variation), confirming robustness to
hyperparameter misspecification.

\subsection{Statistical Significance}

\textbf{Friedman test (coverage).} Across all nine methods and 20 seeds:
$\chi^2 = 3.05$, $p = 0.93$. We fail to reject the null hypothesis of equal
coverage performance.

\textbf{Friedman test (ECE).} $\chi^2 = 34.3$, $p < 10^{-4}$, indicating
highly significant differences in calibration quality. Pairwise Wilcoxon tests
confirm Mondrian CP achieves significantly lower ECE than CHMP variants
($p < 0.01$).

\textbf{Width.} Only CQR ($p = 0.024$, wider) and Ensemble ($p = 0.019$,
narrower) differ significantly from Mondrian CP.

\subsection{Public Real-Data Benchmarks}

Beyond the main synthetic benchmark, we evaluated Mondrian and CHMP on two public
MATPOWER-derived heterogeneous graphs using three seeds each. On ACTIVSg200, both
methods achieved average marginal coverage of $0.935$; CHMP slightly reduced mean
width relative to Mondrian ($0.564$ vs.\ $0.568$) but had slightly higher ECE
($0.073$ vs.\ $0.068$). On IEEE 118, both methods again achieved average marginal
coverage of $0.898$, while CHMP was marginally wider ($0.805$ vs.\ $0.803$) and
less well calibrated by ECE ($0.105$ vs.\ $0.069$). The main conclusion from these
public-network runs is therefore conservative: STRATA transfers across more than
one real topology, but CHMP does not yet demonstrate a consistent out-of-sample
advantage over the simpler Mondrian baseline on these datasets.

\subsection{Alpha Ablation and Fairness Audit}

The alpha ablation confirms the expected conformal tradeoff: average marginal
coverage decreases from $0.974$ at $\alpha=0.05$ to $0.829$ at $\alpha=0.20$,
while average width decreases from $1.166$ to $0.604$. This monotone trend is
important because it demonstrates that STRATA's implementation respects the
intended width--coverage control surface rather than producing unstable behavior
under retuning. A representative group-wise audit at $\alpha=0.10$ produced an
overall coverage of $0.965$ and a worst-group coverage gap of $0.035$ across the
three infrastructure types, suggesting mild but measurable group disparity that
should be monitored in deployment.

% ─────────────────────────────────────────────────────────────
\section{Discussion}
\label{sec:discussion}

\subsection{Width Redistribution vs.\ Aggregate Reduction}

Because all prior graph conformal methods target classification on homogeneous
graphs (\S\ref{sec:related}), STRATA is, to our knowledge, the first
framework to provide distribution-free prediction intervals for node-level
regression on multi-typed graphs. CHMP's primary design goal is not to shrink
intervals in aggregate but to \emph{redistribute} interval width according to
local graph difficulty: nodes in high-residual neighborhoods receive wider
intervals while low-difficulty nodes receive narrower ones. Across 20 synthetic
seeds, CHMP and Mondrian CP produce nearly identical mean widths ($\sim 0.794$),
but the spatial distribution differs meaningfully---confirmed by significant
Moran's I autocorrelation ($p < 0.01$) in the residual field under CHMP. The
real-data comparisons make clear that this redistribution is not automatically
beneficial in aggregate: the width effect is slightly favorable on ACTIVSg200
and slightly unfavorable on IEEE 118. The strongest claim supported by the
current evidence is therefore mechanistic: CHMP changes \emph{where} interval
budget is allocated, which is valuable in infrastructure settings where
spatially non-uniform risk demands non-uniform uncertainty coverage.

\subsection{Limitations}

\begin{enumerate}
\item \textbf{No significant aggregate width improvement.} CHMP does not
produce narrower intervals than Mondrian CP ($p > 0.50$). Only the ensemble
method achieves significant width reduction ($6.4\%$, $p = 0.019$) at
$3\times$ training cost.

\item \textbf{Lambda insensitivity.} Width varies $<0.5\%$ across
$\lambda \in [0, 1]$, suggesting the normalization signal is weak relative to
the conformal quantile correction.

\item \textbf{CQR underperformance.} CQR produces the widest intervals
($0.838$) and lowest coverage ($0.900$).

\item \textbf{ECE trade-off.} CHMP variants exhibit higher ECE than Mondrian
CP ($0.067$--$0.080$ vs.\ $0.057$, $p < 0.01$).

\item \textbf{Synthetic data dominance.} The water and telecom layers in
ACTIVSg200 and IEEE 118 are derived from power demand heuristics rather than
independent utility datasets. This makes the public real-data evaluation useful
for topology transfer, but not a full validation on independently measured
multi-utility networks.

\item \textbf{Mixed real-data evidence.} On the two MATPOWER-based benchmarks,
CHMP matches Mondrian coverage but does not produce a consistent width or ECE
improvement. The method's empirical advantage over simpler baselines therefore
remains benchmark-dependent.

\item \textbf{Scalability.} Current evaluation on graphs with $\sim$400 nodes.

\item \textbf{Static analysis.} STRATA provides static risk intervals;
streaming infrastructure monitoring requires online conformal updates
\citep{gibbs2021adaptive}.

\item \textbf{Residual group disparity.} The fairness audit shows a worst-group
coverage gap of roughly $0.035$ in a representative run. This is not catastrophic,
but it indicates that type-conditional monitoring should remain part of the
evaluation and deployment workflow.
\end{enumerate}

% ─────────────────────────────────────────────────────────────
\section{Conclusion}
\label{sec:conclusion}

We introduced STRATA, the first framework for distribution-free uncertainty
quantification via conformal prediction on heterogeneous graph regression.
While prior graph conformal methods address node \emph{classification} on
homogeneous graphs (DAPS, NAPS, SNAPS, CF-GNN), STRATA targets node-level
\emph{regression} with per-type Mondrian coverage guarantees. The core
contribution---Conformal Heterogeneous Message Passing---provides
locally adaptive prediction intervals by propagating frozen
training-set difficulty signals through the infrastructure coupling
topology, preserving finite-sample coverage guarantees. Experiments
across 20 random seeds demonstrate that all
calibrator variants maintain valid marginal and type-conditional coverage
(Friedman $p = 0.93$).

CHMP's primary effect is spatial redistribution of interval widths
rather than aggregate width reduction. Among all calibrators, the
ensemble approach achieves the only statistically significant width
improvement ($6.4\%$, $p = 0.019$). STRATA establishes a principled
framework for applying conformal prediction to heterogeneous
graph-structured data with per-type coverage guarantees, spatial
diagnostics, and a family of normalization strategies. The current
evidence supports CHMP as a plausible mechanism for difficulty-aware
width redistribution, while also indicating that simpler baselines
remain competitive on public real-data benchmarks.

\acks{This work was conducted independently without external funding.
The author declares no competing interests.
The author consents to publication and confirms that this manuscript is not
under simultaneous review elsewhere.
Code, tests, examples, benchmark scripts, and a supplementary reproducibility
bundle are provided at \url{https://github.com/matthew-lottly/strata}.}

\paragraph{Software.} Code is available at
\url{https://github.com/matthew-lottly/strata}
under the MIT License.
\bibliography{references}

\end{document}
